\section{Consequences for Diophantine solvability}\label{sec:consequences}

For every finite $n\ge1$ and every recursively enumerable
$S\subseteq\NN^n$, Theorem~\ref{thm:main} also gives a finite-fold
representation: each input has at most one complete witness, and therefore
finitely many. The stronger uniqueness assertion yields the promise
undecidability result stated in the introduction. Its proof specializes one
fixed representing polynomial: the degree and number of variables stay fixed,
while the coefficients vary with the input.

\begin{proof}[Proof of Corollary~\ref{cor:unique-h10}]
Fix a nonrecursive recursively enumerable set $S\subseteq\NN$.
Theorem~\ref{thm:main}, with one parameter, supplies a fixed $m\ge1$ and
a fixed polynomial $P\in\ZZ[A,W_1,\ldots,W_m]$ representing $S$
single-fold. Choose one such $P$, put $D=\max\{1,\deg P\}$, and for
each $a\in\NN$ define
\[
 F_a(W_1,\ldots,W_m)=P(a,W_1,\ldots,W_m).
\]
Every $F_a$ has at most one zero in $\NN^m$, and it has a zero exactly
when $a\in S$. This counts the entire remaining tuple
$(W_1,\ldots,W_m)$, including every subsidiary arithmetic variable.
Specialization introduces no variables and does not increase total degree,
so every $F_a$ satisfies the fixed bounds.

Substitution is computable from the finite coefficient encoding of $P$ and
$a$. Although the proof does not numerically specify $P$, the finite code
of this one fixed polynomial can be hard-coded into a Turing machine. Hence $a\mapsto F_a$ is a computable map of encoded inputs.
A hypothetical algorithm with the stated property, applied to $F_a$, would
therefore decide membership in $S$, a contradiction. Every constructed
input satisfies the promise, so the algorithm's behavior outside it is
irrelevant.
\end{proof}

The bounds in Corollary~\ref{cor:unique-h10} are existential: no numerical
values for them or for the fixed polynomial's coefficients are supplied.
The coefficients of $F_a$, and therefore their input bit lengths, may grow
with $a$. The promise permits both zero and one complete solution; it is not an
exactly-one promise. The corollary makes no claim about analogous uniqueness
promises for solutions over $\ZZ$ or $\QQ$.

